\documentclass[10pt]{article}
\usepackage[margin=1in]{geometry}
\usepackage{amsmath,amssymb}
\usepackage{booktabs}
\usepackage{microtype}
\usepackage[hidelinks]{hyperref}

\newcommand{\code}[1]{\texttt{#1}}
\newcommand{\E}{\mathbb{E}}
\newcommand{\wrapop}{\operatorname{wrap}}

\title{Over-the-Air Carrier Phase Synchronization for Distributed Arrays:\\
Foundations, the Rashid--Nanzer Consensus Approach,\\
and the Present Project's Methods and Measurements}
\author{Chandler Stevenson}
\date{July 29, 2026}

\begin{document}
\maketitle

\begin{abstract}
\noindent Part~I develops the required background from first principles:
the coherent-gain arithmetic that motivates phase synchronization, the
physics and mathematics of oscillator noise, the estimation theory of
phase and frequency measurement from pilots, the identifiability limit of
one-way links, and the structure of a closed synchronization loop. Part~II
gives a complete account of M.~Rashid and J.~A.~Nanzer, ``Frequency and
Phase Synchronization in Distributed Antenna Arrays Based on Consensus
Averaging and Kalman Filtering,'' \emph{IEEE Trans.\ Wireless Commun.}
22(4):2789--2802, 2023: its error model, both algorithms, the derivation
logic of its closed-form residual, its reported results, and its
assumptions. Part~III documents this project's work: the simulator and its
ground-truth grading method, a measured closed-loop error budget with
per-term ablation, reciprocal two-way synchronization including the
$\pi$-ambiguity mechanics, a two-level reimplementation of the paper's
algorithms with a head-to-head comparison that exposes an anti-phase
bistability, a two-tier micro-pilot scheme, and a hybrid one-way/two-way
calibration scheme validated under Doppler. All quoted numbers are
reproducible from \code{simulation.py} and the test suite in this
repository.
\end{abstract}

\tableofcontents

%=====================================================================
\part{Foundations}

\section{Coherent gain and the phase-error requirement}
\label{sec:gain}
Let $N$ transmitters radiate the same tone at carrier frequency $f_c$.
At a far-field point the fields add as phasors, giving power
\begin{equation}
P \;\propto\; \Bigl|\sum_{n=1}^{N} e^{j\phi_n}\Bigr|^2 ,
\end{equation}
where $\phi_n$ contains each transmitter's oscillator phase and its
propagation phase to the observation point. If all $\phi_n$ are equal, $P
\propto N^2$: $N$ times the power of a single transmitter times the array
factor $N$. For $N=2$ with phase difference $\Delta\phi$,
\begin{equation}
P \propto |1+e^{j\Delta\phi}|^2 = 2+2\cos\Delta\phi = 4\cos^2(\Delta\phi/2),
\qquad G \equiv \cos^2(\Delta\phi/2),
\label{eq:gain2}
\end{equation}
where $G$ is the fraction of ideal two-element gain; $G(0)=1$,
$G(\pi)=0$. For $N$ elements with independent zero-mean Gaussian phase
errors of standard deviation $\sigma$, $\E[e^{j\phi_n}]=e^{-\sigma^2/2}$,
so the mean coherent power is approximately $N^2e^{-\sigma^2}$ for large
$N$. Setting $\sigma=18^\circ=0.314$\,rad gives $e^{-\sigma^2}=0.906$;
this is the origin of the $18^\circ$ threshold for retaining $90\%$ of
ideal gain used throughout the distributed-array literature.

The same arithmetic governs coherent reception. When two stations receive
a signal from a passive emitter or scatterer and combine (by summation or
cross-correlation), the source's own phase is common to both receive
chains and cancels; the phase difference between the stations' local
oscillators does not. Cross-station coherence is therefore the entire
synchronization requirement for passive coherent sensing.

Scale of the problem: at $f_c=915$\,MHz the carrier period is 1.093\,ns,
so 1\,ns of relative timing error corresponds to
$360^\circ\times0.915=329^\circ$ of carrier phase. Free-running crystal
oscillators disagree by parts in $10^{6}$ (tens of ppm); at 915\,MHz,
1\,ppm is 915\,Hz of carrier offset, accumulating a full cycle of phase
error every 1.09\,ms. Alignment must be continuously maintained, and for
radios connected only by radio, maintained over the air.

\section{Oscillator noise}
\label{sec:osc}

\subsection{State model and noise types}
Write the oscillator's deviation from nominal as phase $\phi$ (rad) and
angular frequency offset $\omega$ (rad/s). Sampled at interval $T$,
\begin{equation}
\begin{bmatrix}\phi_{k+1}\\ \omega_{k+1}\end{bmatrix}
=\begin{bmatrix}1 & T\\ 0 & 1\end{bmatrix}
\begin{bmatrix}\phi_{k}\\ \omega_{k}\end{bmatrix}+\mathbf{w}_k .
\label{eq:oscmodel}
\end{equation}
The process noise $\mathbf{w}_k$ summarizes several physical mechanisms,
classified by the slope of the single-sideband phase-noise spectrum
$\mathcal{L}(f)\propto f^{-\alpha}$:
\begin{description}
\item[White PM ($\alpha=0$):] each sample of phase carries an independent
jitter $\sim\mathcal{N}(0,\sigma_{\mathrm{wpm}}^2)$. An estimator that
averages $L$ samples reduces its contribution by $1/L$.
\item[White FM ($\alpha=2$):] the instantaneous frequency is white, so the
phase is a random walk. If each sample at rate $f_s$ adds an independent
phase step of standard deviation $\sigma_{\mathrm{pn}}$, then over time
$t$
\begin{equation}
\operatorname{var}\bigl(\phi(t)-\phi(0)\bigr)=\sigma_{\mathrm{pn}}^2 f_s\, t
\;\equiv\; k\,t .
\label{eq:walk}
\end{equation}
A random walk has independent increments: the walk accumulated between
two correction epochs cannot be predicted from anything observed before
the first epoch. Consequently $kT$ is an irreducible phase variance for
any synchronizer that corrects every $T$ seconds, independent of SNR,
estimator, or control law.
\item[Flicker FM ($\alpha=3$):] frequency noise with $1/f$ spectrum. It
has no finite-dimensional exact state-space model; the standard surrogate
is a sum of first-order Gauss--Markov (AR(1)) processes whose correlation
times are log-spaced across the timescales of interest, each contributing
a Lorentzian; the sum approximates a $1/f$ spectrum over that range.
\item[Random-walk FM ($\alpha=4$):] white noise added to frequency each
step (the $\sigma_\omega^2$ entry of $\mathbf{w}_k$); models aging and
temperature ramps.
\end{description}

\subsection{Allan deviation and the numbers used here}
The fractional frequency is $y(t)=\dot\phi(t)/(2\pi f_c)$. The Allan
deviation $\sigma_y(\tau)$ is the standard deviation of the difference of
successive $\tau$-averaged values of $y$, divided by $\sqrt2$; each noise
type produces a characteristic slope ($\tau^{-1}$ for white PM,
$\tau^{-1/2}$ for white FM, $\tau^{0}$ for flicker FM, $\tau^{+1/2}$ for
random-walk FM). A quick estimate for white FM: the $\tau$-average of $y$
is $(\phi(t{+}\tau)-\phi(t))/(2\pi f_c\tau)$, with standard deviation
$\sqrt{k\tau}/(2\pi f_c\tau)=\sqrt{k/\tau}/(2\pi f_c)$. The project's
default oscillator uses $\sigma_{\mathrm{pn}}=2\times10^{-4}$\,rad/sample
at $f_s=1$\,MS/s, so $k=0.04$\,rad$^2$/s and
$\sigma_y(1\,\mathrm{s})\approx0.2/(2\pi\cdot915\times10^6)=3.5\times10^{-11}$
--- between a good TCXO and an OCXO. The other defaults: white PM
5\,mrad/sample; flicker 0.05\,Hz RMS (fractional $5.5\times10^{-11}$
floor); frequency random walk 0.1\,Hz per 50\,ms interval; inter-interval
phase process noise 2\,mrad.

\subsection{One crystal, three errors}
A radio synthesizes its carrier and its ADC/DAC sample clock from the
same reference crystal. If the reference runs fractionally fast by $y$,
then simultaneously: the carrier is high by $f_c y$ (CFO); the sample
clock is fast by $y$ (SFO, in ppm, $10^6 y$); and the receiver's frame
timing drifts by $y f_s T$ samples per interval. At the defaults, an
initial CFO of 1500\,Hz at 915\,MHz means $y=1.64\times10^{-6}$: 1.64\,ppm
of SFO and 0.082 samples of timing drift per 50\,ms interval. These are
one random process observed three ways and must co-drift in any faithful
model; the reference's random walk and flicker move all three together.
A digital NCO correction changes none of them --- it shifts baseband
phase/frequency while the physical crystal keeps drifting --- so a model
must track the physical (correction-free) frequency separately from the
corrected state to derive SFO and timing.

\section{Measuring phase and frequency from pilots}
\label{sec:measure}

\subsection{The phase CRLB}
Let $r_n = e^{j\phi}s_n + w_n$, $n=1{:}L$, with known unit-power pilot
$s_n$ and circular complex Gaussian noise of variance $\sigma_w^2$
($\mathrm{SNR}=1/\sigma_w^2$). The log-likelihood depends on $\phi$
through $\tfrac{2}{\sigma_w^2}\operatorname{Re}\{e^{-j\phi}\sum_n r_n
s_n^*\}$; differentiating twice and taking the expectation gives Fisher
information $I(\phi)=2L\,\mathrm{SNR}$, hence
\begin{equation}
\operatorname{var}(\hat\phi)\;\ge\;\frac{1}{2\,\mathrm{SNR}\,L}.
\label{eq:crlbphase}
\end{equation}
At 20\,dB SNR ($\mathrm{SNR}=100$) with $L=2\times2047$ pilot samples,
the bound is $1.22\times10^{-6}$\,rad$^2$, i.e., 1.1\,mrad.

\subsection{Frequency from two pilot blocks}
Estimate the phase of two identical pilot blocks of $L_b$ samples whose
centers are separated by $T_{\mathrm{sep}}$, and form
$\hat\omega=(\hat\phi_2-\hat\phi_1)/T_{\mathrm{sep}}$. Then
\begin{equation}
\operatorname{var}(\hat\omega)=\frac{2}{2\,\mathrm{SNR}\,L_b\,T_{\mathrm{sep}}^2}
=\frac{1}{\mathrm{SNR}\,L_b\,T_{\mathrm{sep}}^2}\ \ \text{(rad/s)}^2 .
\label{eq:crlbfreq}
\end{equation}
With $L_b=2047$, $T_{\mathrm{sep}}=2175$\,samples\,$=2.175$\,ms at
1\,MS/s, and 20\,dB SNR: $\sigma_{\hat\omega}=1.02$\,rad/s, i.e.,
0.16\,Hz.

\subsection{Phase noise inside the observation}
The receiver's LO phase-noise walk runs during the pilot itself. The
phase estimate is approximately the average of the walk over the pilot
span. For a walk $W_i$ with per-sample variance $\sigma_{\mathrm{pn}}^2$
averaged over samples $[a, a{+}n]$,
$\operatorname{var}(\overline{W})\approx\sigma_{\mathrm{pn}}^2(a+n/3)$
(the first term because the walk enters the window already displaced, the
second from averaging a Brownian path). At the default frame (long fields
spanning roughly samples 256--4606),
$\sigma_{\mathrm{pn}}^2(256+4350/3)\approx6.8\times10^{-5}$\,rad$^2$:
8.3\,mrad, seven times the AWGN bound \eqref{eq:crlbphase}. The same walk
decorrelates the two blocks of the frequency estimator, adding variance
$\sigma_{\mathrm{pn}}^2 N_{\mathrm{sep}}/T_{\mathrm{sep}}^2\approx18$
(rad/s)$^2$: 0.68\,Hz, four times the AWGN frequency bound. At these
settings estimation error is phase-noise-limited, not SNR-limited, and a
filter designed from the oscillator's data sheet must include these terms
in its measurement covariance or it will overweight its measurements.

\subsection{Waveforms: Zadoff--Chu sequences and the detection metric}
A Zadoff--Chu (ZC) sequence of odd length $L$ and root $u$ coprime to $L$
is $x[n]=\exp(-j\pi u\,n(n{+}1)/L)$. It has constant amplitude (bounded
PAPR through the PA) and ideal periodic autocorrelation (zero at all
nonzero cyclic shifts), so a matched filter produces a single sharp
timing/phase peak even in multipath, and a cyclic prefix converts linear
convolution with the channel into cyclic, preserving the property. The
project's frame is: a short training field (STF), a length-16 ZC repeated
16 times (256 samples); then two long training fields (LTF), each a
128-sample cyclic prefix plus a length-2047 ZC (root 25), total 4606
samples ($=4.6$\,ms at 1\,MS/s).

Detection and coarse CFO use the STF's periodicity. With lag $D=16$,
\begin{equation}
P(d)=\sum_{m} r^*[d{+}m]\,r[d{+}m{+}D],\qquad
M(d)=\frac{|P(d)|}{\sqrt{\textstyle\sum|r[d{+}m]|^2\sum|r[d{+}m{+}D]|^2}} ,
\end{equation}
maximized over $d$; $M$ near 1 indicates a repeated structure regardless
of channel or CFO. The angle of $P$ estimates CFO:
$\hat\omega=\angle P/(D\,T_s)$, unambiguous while $|\omega|D T_s<\pi$,
i.e., $|f|<f_s/(2D)=31.25$\,kHz --- this sets the acquisition range.
Fine timing then maximizes the normalized cross-correlation against the
whole known preamble; fine CFO comes from the phase of the correlation
between the two LTFs (separation 2175 samples,
\eqref{eq:crlbfreq}); the frame phase comes from matched-filtering the
LTFs after CFO correction, referenced to the frame center.

\subsection{The identifiability limit of one-way measurement}
The received pilot phase is
\begin{equation}
m \;=\; \underbrace{(\phi_{\mathrm{tx}}-\phi_{\mathrm{rx}})}_{\theta}
\;+\; \phi_c \;+\; n,
\label{eq:oneway}
\end{equation}
where $\phi_c$ is the channel phase (propagation delay modulo the carrier
period, plus the multipath composite). $\theta$ and $\phi_c$ enter
identically; no processing of one-way data separates them. A one-way
synchronizer that drives $m\to0$ therefore forces the receiving
oscillator to absorb $\phi_c$ (a constant $+2.92$\,rad at the project's
default channel realization). Whether this matters depends on the
application:
\begin{itemize}
\item If a downstream channel-estimation step exists (a user measuring
pilots from both stations, as in cell-free MIMO), any static per-station
phase is absorbed into the user's channel estimate; only the
\emph{drift} of the differential phase between estimates matters. One-way
synchronization suffices.
\item With no feedback path (open-loop beamforming at a computed
direction, passive coherent sensing), the carriers must align at the
antennas, and the $\phi_c$ bias is disqualifying.
\end{itemize}
The remedy is reciprocity. Exchanging pilots in both directions over the
same channel within a coherence time, with
$m_{\mathrm{f}}=\theta+\phi_c+n_1$ and $m_{\mathrm{r}}=-\theta+\phi_c+n_2$,
\begin{equation}
\tfrac12(m_{\mathrm{f}}-m_{\mathrm{r}})=\theta+\tfrac12(n_1{-}n_2),\qquad
\tfrac12(m_{\mathrm{f}}+m_{\mathrm{r}})=\phi_c+\tfrac12(n_1{+}n_2),
\label{eq:twoway}
\end{equation}
and the two combinations are uncorrelated
($\operatorname{cov}=\tfrac14(\operatorname{var}n_1-\operatorname{var}n_2)=0$
for equal noise). Each has half the one-way noise variance. Two caveats
carry through the rest of this document. First, phases are measured
modulo $2\pi$, so the half-difference is defined modulo $\pi$: if
$(m_{\mathrm{f}}-m_{\mathrm{r}})$ is computed wrapped, adding $2\pi$
before halving shifts the result by $\pi$. The branch must be chosen, and
a wrong choice is self-consistent (Section~\ref{sec:pi}). Second, only
the propagation channel is reciprocal; each node's transmit and receive
chains differ. With chain phases $\tau^{\mathrm{tx}}_i,
\tau^{\mathrm{rx}}_i$, the forward measurement contains
$\tau^{\mathrm{tx}}_A+\tau^{\mathrm{rx}}_B$ and the reverse
$\tau^{\mathrm{tx}}_B+\tau^{\mathrm{rx}}_A$, so the half-difference
carries a bias
$\tfrac12[(\tau^{\mathrm{tx}}_A-\tau^{\mathrm{rx}}_A)-
(\tau^{\mathrm{tx}}_B-\tau^{\mathrm{rx}}_B)]$
that does not cancel and must come from loopback calibration.

\section{The closed loop}
\label{sec:loop}

\subsection{EKF tracking of phase and frequency}
The tracked state is the relative offset $x=[\theta,\omega]^\top$ with
dynamics \eqref{eq:oscmodel}. To avoid wrap discontinuities the phase is
observed through its cosine and sine: the measurement vector is
$z=[\cos m,\ \sin m,\ \hat\omega_{\mathrm{meas}}]^\top$ with model
$h(x)=[\cos\theta,\sin\theta,\omega]^\top$ and Jacobian
\begin{equation}
H(x)=\begin{bmatrix}-\sin\theta & 0\\ \cos\theta & 0\\ 0 & 1\end{bmatrix}.
\end{equation}
Predict: $\hat x\leftarrow F\hat x$, $P\leftarrow FPF^\top+Q$. Update
(iterated, because $\sin$ is far from linear at large innovations):
starting from the prior $(\hat x^-,P^-)$, repeat
\begin{equation}
K=P^-H^\top(HP^-H^\top+R)^{-1},\qquad
\hat x^{(i+1)}=\hat x^- + K\bigl[z-h(\hat x^{(i)})-H(\hat x^--\hat x^{(i)})\bigr]
\end{equation}
with $H$ re-linearized at $\hat x^{(i)}$, until convergence (six
iterations suffice here); then the Joseph-form covariance update
$P=(I{-}KH)P^-(I{-}KH)^\top+KRK^\top$, which preserves symmetry and
positive semidefiniteness. $Q$ is assembled from
Section~\ref{sec:osc} (both nodes' random walks, the white-FM walk over
one interval, the flicker innovation); $R$ from
Section~\ref{sec:measure} (AWGN CRLB, intra-observation phase noise,
white PM).

\subsection{Actuation, latency, and the resulting error term}
The correction is applied by a numerically controlled oscillator with
finite resolution (here: phase quantized to $2\pi/2^{16}$, frequency to
0.01\,Hz) and finite processing latency: a command computed from the
frame at epoch $k$ loads at epoch $k{+}\ell$. A controller that knows
$\ell$ commands the \emph{predicted} state $F^{\ell}\hat x$, which
cancels the deterministic drift $\omega\ell T$; what cannot be cancelled
is (i) the process noise that occurs during the latency (the walk term,
already counted per interval) and (ii) the error in the frequency
estimate itself, which integrates into phase as
\begin{equation}
\sigma_{\theta,\mathrm{latency}}\;=\;\sigma_{\hat\omega}\,\ell\,T ,
\label{eq:latency}
\end{equation}
where $\sigma_{\hat\omega}$ is the filter's steady-state frequency
posterior. That posterior follows from the scalar Riccati recursion
$P^-=P+Q_\omega$, $P=(P^-R_\omega)/(P^-+R_\omega)$: with
$Q_\omega=2(2\pi\cdot0.1)^2=0.79$ and $R_\omega\approx1.03$ (rad/s)$^2$
(the AWGN and phase-noise terms above), the fixed point is $P=0.59$,
$\sigma_{\hat\omega}=0.77$\,rad/s, so
$\sigma_{\theta,\mathrm{latency}}\approx0.77\times0.05=38$\,mrad at
$\ell=1$, $T=50$\,ms. Together with the walk term
$\sqrt{kT}=45$\,mrad this predicts the measured floors of
Part~III before any simulation is run.

%=====================================================================
\part{The Reference Paper: Rashid \& Nanzer, IEEE TWC 2023}

\section{Setting and error model}
The paper treats $N$ nodes forming an open-loop coherent distributed
array: nodes coordinate mutually, with no feedback from the beamforming
destination. Nodes form an undirected graph $\mathcal{G}=(\mathcal{V},
\mathcal{E})$. Each node broadcasts its oscillator signal
$s_n(t)=e^{j(2\pi f_n t+\theta_n)}$; neighbors estimate $(f_n,\theta_n)$
from $L=Tf_s$ received samples. The total phase error at node $n$ (their
Eq.~1) is
\begin{equation}
\delta\phi_n = 2\pi\,\delta\!f_n T + 2\pi\varepsilon_f T + \theta^e_n
+ \delta\theta^f_n + \delta\theta_n + \varepsilon_\theta ,
\label{eq:their1}
\end{equation}
respectively: phase accrued by the frequency offset over the update
interval; by the frequency estimation error; a constant
$\theta^e_n\sim\mathcal{U}[0,2\pi)$ lumping geometry, timing
misalignment, and channel-induced phase (``assumed corrected'');
intra-interval drift; phase jitter; and phase estimation error. The
statistics:
\begin{itemize}
\item Frequency drift: $\sigma_f=f_c\sqrt{\beta_1/T+\beta_2T}$ (white FM
plus random-walk FM in Allan-deviation form), with
$\beta_1=\beta_2=5\times10^{-19}$ for a quartz TCXO.
\item Intra-interval drift: $\sigma^p_\theta=\pi T\sigma_f$, from
integrating a linearly-varying frequency offset across the interval
(their Eqs.~3--4, 13).
\item Phase jitter: $\sigma_\theta=\sqrt{2\cdot10^{A/10}}$ with
integrated phase-noise power $A=-53.46$\,dBc, giving 2.7\,mrad ---
\emph{drawn independently every interval}.
\item Estimation errors at the CRLBs: $\sigma^m_f=\sqrt{6/((2\pi)^2L^3
\mathrm{SNR})}$ (normalized frequency; their Eq.~16) and
$\sigma^m_\theta=\sqrt{2/(L\,\mathrm{SNR})}$ (their Eq.~17). Under TDMA
with $D$ average neighbors, $L\to L/D$, inflating both.
\end{itemize}

\section{DFPC: decentralized frequency and phase consensus}
Stack node frequencies into $\mathbf f$ and phases into
$\boldsymbol\theta$. One iteration per interval $T$:
\begin{equation}
\mathbf f(k)=\mathbf W\,\mathbf f(k{-}1),\qquad
\boldsymbol\theta(k)=\mathbf W\,\boldsymbol\theta(k{-}1),
\label{eq:consensus}
\end{equation}
with the Metropolis--Hastings mixing matrix
$w_{ij}=1/(1+\max\{n_i,n_j\})$ for $(i,j)\in\mathcal E$,
$w_{ii}=1-\sum_{j\ne i}w_{ij}$, $w_{ij}=0$ otherwise ($n_i$ = degree of
node $i$). $\mathbf W$ is symmetric doubly stochastic; its largest
eigenvalue is 1 with eigenvector $\mathbf 1$, so $\mathbf W^{I}\to
\tfrac1N\mathbf{11}^\top$ and every node converges to the arithmetic
mean of the initial states. Convergence rate and error accumulation are
governed by the second-largest eigenvalue modulus $\lambda_2$: sparse
graphs have $\lambda_2$ near 1 (slow mixing), dense graphs near 0. Row
$n$ of \eqref{eq:consensus} involves only node $n$'s neighbors, so the
computation is local; and because the broadcast \emph{waveform} carries
the state, plain DFPC needs no data link --- neighbors estimate the
state, which is why the estimation errors are injected into the
recursion (their Algorithm~1 adds the drift, jitter, and estimation
draws each iteration before mixing).

\subsection{Derivation of the steady-state residual (their Eq.~27)}
With per-iteration error vector $\mathbf e_k$ of per-node variance
$\sigma_e^2$, the state after $I$ iterations is
\begin{equation}
\mathbf z(I)=\mathbf W^{I}\mathbf z(0)+\sum_{i=0}^{I-1}\mathbf W^{i+1}
\mathbf e_{I-i}.
\end{equation}
The first term converges to the average of initial values; the error sum,
measured relative to the network average (subtract
$\tfrac1N\mathbf{11}^\top$), has covariance
$\sigma_e^2\sum_{m=1}^{I}(\mathbf W-\tfrac1N\mathbf{11}^\top)^{2m}$,
bounded by $\sigma_e^2\sum_m\lambda_2^{2m}$. For sparse graphs
($\lambda_2\to1$) the geometric sum is large and the residual saturates
at the injected level itself; the total phase spread is then the
root-sum-of-squares of the five per-iteration error terms:
\begin{equation}
\sigma_{\phi,\mathrm{total}}
=\sqrt{(2\pi\sigma_fT)^2+(\sigma^m_\phi)^2+(\pi T\sigma_f)^2
+(\sigma^m_\theta)^2+\sigma_\theta^2}.
\label{eq:eq27}
\end{equation}
Dense graphs mix errors away faster and sit below this bound. Equation
\eqref{eq:eq27} is the paper's central analytical result: a closed-form
floor in oscillator data-sheet quantities, CRLBs, and the update
interval. Note what is absent: no actuation latency (updates are
instantaneous state assignments), no channel term (swallowed by
$\theta^e_n$), and the jitter re-drawn per interval has no memory.

\section{KF-DFPC: per-node Kalman filtering before averaging}
At small $T$ the CRLB terms dominate \eqref{eq:eq27} (fewer samples per
estimate). The paper's remedy: each node runs a two-state Kalman filter
on its own noisy state observation, then the network consensus-averages
the \emph{posterior means}. State $[f_n,\theta_n]$, process covariance
(their Eq.~29)
\begin{equation}
\mathbf Q=\begin{bmatrix}\sigma_f^2 & -\pi T\sigma_f^2\\
-\pi T\sigma_f^2 & \pi^2T^2\sigma_f^2+\sigma_\theta^2\end{bmatrix},
\end{equation}
measurement covariance $\boldsymbol\Sigma=
\mathrm{diag}((\sigma^m_f)^2,(\sigma^m_\theta)^2)$, standard
predict/update (their Eqs.~33, 35--36). After mixing the means with
$\mathbf W$, the per-node covariances are re-initialized by mixing their
entries with squared weights (their Eq.~39). Per-node complexity is
$\mathcal O(\mathrm{card}\{\chi_n\}+8)$ per iteration. The variant
requires each node to broadcast its posterior mean \emph{and covariance}
to neighbors --- an explicit data side channel, assumed error-free and
instantaneous; plain DFPC does not need one.

Reported results (their Figs.~5--11), at $f_c=1$\,GHz, $f_s=10$\,MHz,
$T=0.1$\,ms, SNR $=0$\,dB, $10^3$ trials: KF-DFPC converges in far fewer
iterations on sparse graphs (2 vs.\ 16 at $N{=}100$, connectivity 0.05;
4 vs.\ 158 at $N{=}65$); its converged error is nearly independent of SNR
(the filter compensates measurement quality with memory); it outperforms
diffusion LMS, diffusion KF, and the Kalman consensus information filter
at low SNR and sparse connectivity; and sweeping $T$ shows a DFPC optimum
near $T=20$\,ms --- drift errors grow with $T$, estimation errors shrink
with $T$ --- while KF-DFPC holds the theoretical floor for
$T\lesssim20$\,ms. No closed form is derived for KF-DFPC's steady state;
it is characterized by simulation against \eqref{eq:eq27}.

\section{The paper's assumptions, stated precisely}
\label{sec:assess}
(i) The measurement model (their Eq.~15) is
$x_n(t)=e^{j(2\pi f_nt+\theta_n)}+n_o(t)$: each node's oscillator is
observed \emph{directly} in AWGN. There is no propagation channel in the
measurement, so the identifiability problem of
Section~\ref{sec:measure} cannot arise, and all channel/geometry phase
sits in the uniform constant $\theta^e_n$ assumed corrected elsewhere.
(ii) Consensus updates are instantaneous assignments; no latency, NCO
quantization, or control dynamics. (iii) Jitter is i.i.d.\ per interval;
intra- and inter-interval phase noise are statistically disconnected;
flicker FM is absent from the analytical model. (iv) Validation is
statistics-level Monte Carlo: their Algorithm~1 injects Gaussian draws
from the assumed distributions, so simulations test the algebra, not the
physics. (v) Time synchronization is assumed complete; the KF side
channel is ideal. These are scope choices, not errors; Part~III measures
their consequences.

%=====================================================================
\part{This Project}

\section{The simulator}
\label{sec:sim}
The \code{ota\_sync} package simulates the link at sampled-IQ level.
Defaults:

\begin{center}
\begin{tabular}{llll}
\toprule
carrier & 915\,MHz & sample rate & 1\,MS/s\\
sync interval $T$ & 50\,ms & SNR (nominal) & 20\,dB\\
channel & 3GPP TDL-D, 100\,ns DS & initial CFO & 1500\,Hz\\
shadowing & 2\,dB, AR(1), 10\,s & white FM & $2{\times}10^{-4}$\,rad/sample\\
white PM & 5\,mrad/sample & flicker FM & 0.05\,Hz RMS\\
freq.\ random walk & 0.1\,Hz/interval & ADC/DAC & 12\,bit\\
NCO quantization & $2\pi/2^{16}$, 0.01\,Hz & correction latency & 1 interval\\
\bottomrule
\end{tabular}
\end{center}

Per interval, the transmit chain is: preamble (Section
\ref{sec:measure}) modulated on the master carrier; soft PA limiter; DAC
quantization; Sionna TDL channel taps (per-interval snapshots,
Jakes-correlated across intervals when mobility is nonzero); AR(1)
log-normal shadowing gain; additive noise at a \emph{fixed} thermal floor
established once at nominal channel gain, so per-frame SNR fades with
the channel (deriving noise from each frame's received power, a common
shortcut, silently pins SNR and removes fading). The receive chain:
resampling at the current physical SFO with the accumulated fractional
timing offset (whole-sample drift steps the arrival inside the capture
window --- a receiver re-centering on each detection --- and the
sub-sample residual is a fractional delay); downconversion by the slave
LO; the intra-frame white-FM walk (whose end value is folded back into
the oscillator state, plus an independent increment for the dead time,
making intra- and inter-interval phase noise one continuous process);
white-PM jitter; IQ gain/phase imbalance; AGC; DC offset; ADC
quantization. The receiver DSP is the STF/LTF processing of
Section~\ref{sec:measure}; the EKF of Section~\ref{sec:loop} tracks;
corrections are forward-predicted by the known latency, quantized, and
queued.

\subsection{Ground-truth grading: the oracle path}
\label{sec:oracle}
A grading pitfall: an early metric compared the EKF posterior to the
measurement that had just updated it. With Kalman gain $\kappa$, that
residual is $(1-\kappa)$ times the innovation; at $\kappa\approx0.9$ it
understates the true error by an order of magnitude, and the code
reported 0.105\,mrad tracking RMSE. The corrected method builds, for
every capture, an \emph{oracle twin} sharing everything deterministic ---
frame, timing jitter draw, channel realization, shadowing state, PA/DAC
distortion, sample-clock timing --- and omitting everything stochastic or
receiver-imposed: AWGN, phase noise, IQ imbalance, DC offset, ADC. The
same estimator run on the oracle capture defines the true observable
phase and frequency; every error metric is computed against these. The
tracking RMSE under this grading is 12.0--12.5\,mrad, a factor
$\sim$100 above the self-referential figure. Distinct from the paper's
validation (Section~\ref{sec:assess}, iv): here the closed-form claims
are checked against an independent physical model rather than against
their own assumptions.

\section{Measured closed-loop error budget}
\label{sec:budget}
Combining Sections \ref{sec:osc}--\ref{sec:loop}, the steady-state
closed-loop phase residual of the one-way loop is
\begin{equation}
\sigma^2 \approx \underbrace{\sigma_{\mathrm{pn}}^2 f_s T}_{\text{white-FM
walk per interval}}
+\underbrace{(\sigma_{\hat\omega}\,\ell T)^2}_{\text{frequency posterior
$\times$ latency}}
+\underbrace{\sigma^2_{\mathrm{track}}}_{\text{measurement/tracking}} .
\label{eq:budget}
\end{equation}
Predicted terms at defaults: 45\,mrad, 38\,mrad
(Section~\ref{sec:loop}), and $\sim$12\,mrad. Ablation results (100
intervals):
\begin{center}
\begin{tabular}{lcc}
\toprule
Configuration & Tracking RMSE & Closed-loop residual RMS\\
\midrule
Full defaults & 12.0\,mrad & 69.6\,mrad\\
RF impairments off (isolates the latency term) & 1.1\,mrad & 34.3\,mrad\\
Zero latency (isolates the tracking term) & 12.3\,mrad & 12.3\,mrad\\
\bottomrule
\end{tabular}
\end{center}
The impairment-free case measures the latency term alone at 34.3\,mrad
against the 38\,mrad Riccati prediction; the impairment-free tracking
RMSE of 1.1\,mrad sits on the AWGN CRLB \eqref{eq:crlbphase}. The
full-default residual exceeds the root-sum-of-squares of the three
predicted terms (60\,mrad) by the flicker and shadowing contributions not
in \eqref{eq:budget}. Consequences: above moderate SNR the floor is set
by oscillator quality, correction cadence, and latency, not pilot energy;
the levers are a quieter reference, a shorter interval, or lower latency.
Relative to \eqref{eq:eq27}: the latency term \eqref{eq:latency} has no
counterpart there (updates are instantaneous); the Kalman steady-state
posterior replaces per-measurement CRLBs; flicker and process continuity
are added. A literature pass over 19 primary sources (each claim verified
by three-referee adversarial review against full text) found no published
closed-loop residual containing the latency-coupled term through
December 2025; the closest works are Eq.~\eqref{eq:eq27} itself and
open-loop drift-versus-update-interval analyses (Mghabghab et al., IEEE
Access 2021).

\section{Two-way synchronization and the $\pi$ ambiguity}
\label{sec:pi}
Implementing \eqref{eq:twoway} in the simulator: forward and reverse
frames traverse the same channel-tap realization and the same shadowing
state within the interval (reciprocity), with the intra-frame LO walk of
the forward frame folded into the oscillator state before the reverse
frame so the two directions see one continuous noise process. The EKF
consumes the half-difference phase and half-difference frequency
(channel Doppler is common to both directions and cancels from the
frequency too), with half the one-way measurement covariance. Result:
the raw oscillator residual --- unresolvable at $-2.92$\,rad under
one-way sync --- becomes \textbf{81.8\,mrad RMS} (two-station gain
99.83\%) at 19.1\% airtime.

The modulo-$\pi$ branch of the half-difference is chosen at acquisition
(candidate nearest zero) and thereafter the candidate nearest the filter
prediction. A wrong acquisition branch produces a lock at relative phase
$\pi$ in which every subsequent measurement is self-consistent: the
filter tracks its shifted belief exactly, and nothing internal to the
loop reveals the error. This project first encountered the failure only
after many correct runs, for an arithmetic reason worth recording: the
default CFO times the interval, $1500\times0.05=75.000$ carrier cycles,
is an exact integer, so the initial 1.2\,rad offset ($<\pi/2$)
reappeared unchanged at every acquisition and the nearest-zero rule
happened to be right. A later configuration with 5\,ms sub-intervals
produced 7.5 cycles per step, the acquisition-time phase landed at
1.94\,rad ($>\pi/2$), and the loop locked at $\pi$ with 0.01\% coherent
gain. The deployed resolution is a one-time calibration after lock: any
coarse external power measurement tells constructive from destructive
(one bit); if destructive, the NCO is flipped by $\pi$. Implementation
detail with a physical meaning: the filter must \emph{not} be reset at
the flip. The half-difference measurement is invariant under $\pi$
shifts of the true phase, so the filter's near-zero state simply
re-attaches to the true branch; resetting it by $\pi$ re-creates the
shifted self-consistent lock (this was observed before being understood).

\section{Reimplementation of DFPC/KF-DFPC and the comparison}
\label{sec:compare}
The paper's algorithms were implemented at two levels.
\emph{Statistics level}, faithful to Algorithms 1--2 including the
Metropolis--Hastings weights, the graph sampler, and the Eq.~39
covariance mixing: at their canonical setup ($N{=}20$, connectivity 0.2,
$T{=}0.1$\,ms, 0\,dB) DFPC converges to $8.0^\circ$ against the
\eqref{eq:eq27} bound of $5.9^\circ$ (consistent: moderate connectivity
sits near, sparse below the bound is not expected), and KF-DFPC reduces
it to $2.1^\circ$ --- both of the paper's qualitative claims reproduce.
\emph{Physical level}, at $N{=}2$ over the full simulator: the two-node
Metropolis--Hastings weight is $w_{12}=\tfrac12$, so each node retunes
its own NCO by half its measured offset toward the other, with the same
latency, quantization, and calibration machinery as the project's own
loops.

\subsection{Anti-phase capture of naive over-the-air consensus}
Run as published --- each node consensus-averages its \emph{raw one-way}
measurement, consistent with the paper's assumption that channel phase
is handled elsewhere --- the update in the relative coordinate
$r=\theta_A-\theta_B$ is
\begin{equation}
r' \;=\; r+\tfrac12\wrapop(-r+\phi_c)-\tfrac12\wrapop(r+\phi_c),
\label{eq:bistable}
\end{equation}
since node $A$ adds half its measurement $\wrapop(-r+\phi_c)$ to itself
and node $B$ adds half of $\wrapop(r+\phi_c)$. Fixed-point analysis:
\begin{itemize}
\item If neither argument wraps, the wraps are identities and
$r'=r-r=0$: the reciprocal channel phase enters both nodes' updates
identically and cancels; the pair aligns in one step regardless of
$\phi_c$.
\item $r=\pi$ is also a fixed point: $\wrapop(-\pi+\phi_c)
=\wrapop(\pi+\phi_c)$ because $-\pi\equiv\pi$, so the two halves cancel
and $r'=\pi$. Perturbations decay back (the linearization about $\pi$
matches that about 0), so anti-phase is \emph{stable}.
\item If exactly one argument crosses $\pm\pi$ during convergence ---
which happens when $|r|+|\phi_c|>\pi$ --- the two wraps disagree by
$2\pi$ and the update lands on the $\pi$ branch.
\end{itemize}
Worked example at the default realization ($\phi_c=+2.922$, $r_0=1.2$):
$\wrapop(1.2+2.92)=-2.16$, $\wrapop(-1.2+2.92)=+1.72$,
$r_1=1.2+\tfrac{1.72}{2}+\tfrac{2.16}{2}=3.1416=\pi$. Verification was
threefold: the full simulation reproduces $r_1=-3.117$ (the same point on
the circle, 25\,mrad of measurement noise from $\pi$), byte-identically
across reruns; a four-line iteration of \eqref{eq:bistable} with no radio
code shows capture for $\phi_c=\pm2.92$ and alignment for
$\phi_c=-1.0,+0.5$; and across seeds, capture occurs exactly when
$|r_0|+|\phi_c|>\pi$ predicts (channel phase $+2.92$: capture;
$-1.50,+1.37,-1.39$: alignment). The failure is invisible to the paper's
channel-free model and is realization-dependent: the same hardware and
algorithm aligns on one channel and locks anti-phase on another.

\subsection{Head-to-head results, identical physical conditions}
\begin{center}
\begin{tabular}{lccc}
\toprule
Approach & Residual RMS & 2-station gain & Airtime\\
\midrule
Two-tier micro-pilot (\S\ref{sec:micro}) & 28.1\,mrad & 99.98\% & 26.0\%\\
Hybrid one-way/anchors (\S\ref{sec:hybrid}) & 33.8\,mrad & 99.97\% & 14.9\%\\
Two-way EKF (\S\ref{sec:pi}) & 80.8\,mrad & 99.84\% & 19.1\%\\
KF-DFPC + reciprocity (paper, steelman) & 82.8\,mrad & 99.83\% & 19.1\%\\
DFPC + reciprocity (paper, unfiltered) & 153.0\,mrad & 99.42\% & 19.1\%\\
DFPC naive (paper, as published) & 3009\,mrad & 0.68\% & 19.1\%\\
\bottomrule
\end{tabular}
\end{center}
The ``+ reciprocity'' rows exchange measurements over the side channel
the paper already assumes for KF-DFPC and consense on the channel-free
half-difference. Three findings, each independently checked:
\begin{enumerate}
\item Unfiltered consensus passes raw frequency measurement noise
straight into the NCOs each interval. The measured steady frequency
residual is 0.442\,Hz; the implied drift per 50\,ms interval is
$2\pi\times0.442\times0.05=139$\,mrad, which with the 45\,mrad walk
predicts $\sqrt{139^2+45^2}=146$\,mrad against 153 observed. The paper's
claim that filtering matters at short intervals reproduces over a
physical link.
\item Filtered consensus and the project's two-way EKF land on the same
floor (82.8 vs.\ 80.8\,mrad; within a few mrad at every tested seed; not
bit-identical, ruling out shared state). Both are Kalman-filtered
reciprocal loops differing only in where the correction is applied. At
fixed physics, the algorithm layer is interchangeable; the floor is
\eqref{eq:budget}.
\item The naive variant's collapse (0.68\% of ideal gain) is the
measured cost of the channel-free measurement assumption.
\end{enumerate}

\section{Two-tier micro-pilot synchronization}
\label{sec:micro}
Both dominant terms of \eqref{eq:budget} scale with the correction
period, and phase --- unlike detection, timing, or wide-range CFO ---
does not need a long pilot once the loop is locked: a short known
sequence, correlated at the expected arrival after derotating by the
current frequency estimate, gives a phase measurement with variance
$1/(2\,\mathrm{SNR}\,L_\mu)$ plus the (much smaller) intra-pilot phase
noise. The two-tier scheme keeps the full two-way frame once per 50\,ms
interval (detection, timing, CFO, and the anchor for acquisition) and
inserts $M$ reciprocal phase-only micro-pilots --- a 32-sample cyclic
prefix plus a ZC-255, 287 samples $=0.29$\,ms per direction --- at evenly
spaced sub-intervals, with corrections issued every sub-interval.
Tracking-mode feasibility at the defaults: timing drift per 10\,ms
sub-interval is 0.016 samples (search $\pm4$ suffices); derotation error
from a 0.2\,Hz residual CFO over 287\,$\mu$s is 0.36\,mrad. Measured
residual--airtime trade:
\begin{center}
\begin{tabular}{lcccc}
\toprule
Micro-pilots per interval $M$ & 0 & 2 & 4 & 9\\
\midrule
Residual RMS (mrad) & 71.5 & 35.3 & 27.9 & 22.8\\
Airtime & 19.1\% & 22.6\% & 26.0\% & 34.6\%\\
\bottomrule
\end{tabular}
\end{center}
The scaling follows \eqref{eq:budget} with $T\to T/(M{+}1)$: at $M=4$
the predicted walk term is $2\times10^{-4}\sqrt{10^4}=20$\,mrad and the
latency term $0.77\times0.01=8$\,mrad, so $\sqrt{20^2+8^2+\text{track}^2}
\approx23$--25\,mrad against 27.9 measured (the full-frame substep and
flicker account for the difference). Returns diminish beyond $M\approx4$
as the measurement and staleness terms take over. It was the $M=9$
configuration (5\,ms substeps, 7.5 carrier cycles per step) that exposed
the acquisition-branch failure of Section~\ref{sec:pi}.

\section{Hybrid one-way/two-way calibration}
\label{sec:hybrid}
Reciprocal exchange costs both directions' airtime, and
\eqref{eq:twoway} shows what the second direction buys: separation of
$\theta$ from $\phi_c$. For fixed stations, $\phi_c$ is nearly constant
(the TDL taps are static at zero speed; shadowing changes amplitude, not
phase) while $\theta$ wanders at the oscillator rate. The processes have
disjoint spectral support, so a joint estimator can track both with
mostly one-way data. State $x=[\theta,\omega,\phi_c]^\top$; dynamics
diagonal except the $\theta\!-\!\omega$ integrator; process noise
$q_\theta$ (walk + oscillator), $q_\omega$ (random walk + flicker
innovation), and $q_c=\sigma_c^2$, the assumed channel-phase drift per
interval --- the estimator's statement of channel coherence time.

Observations:
\begin{itemize}
\item \textbf{One-way pilots} (a full frame each interval for
timing/CFO; one-way micro-pilots at sub-intervals): phase observation
$\theta+\phi_c$, Jacobian row $[-\sin(\cdot),\,0,\,-\sin(\cdot)]$ ---
identical sensitivity to $\theta$ and $\phi_c$. The direction
$\theta-\phi_c$ is unobservable from one-way data alone; between
anchors, the split is maintained only by the process priors (innovations
are attributed in proportion to $q_\theta$ vs.\ $q_c$).
\item \textbf{Two-way anchors} every $K$ intervals: the half-difference
observes $\theta$ (branch-resolved against the filter), the
half-difference of directional frequencies observes $\omega$ free of
channel Doppler, and $\phi_c$ is recovered as
$m_{\mathrm f}-\hat\theta$, which algebraically equals the half-sum
\eqref{eq:twoway} and is uncorrelated with the half-difference.
\end{itemize}
NCO corrections act on $[\theta,\omega]$ only (the channel cannot be
actuated). At a $\pi$-calibration flip, the true $\theta$ moves by
$\pi$ while the one-way observable $\theta+\phi_c$ must remain matched to
the filter's; the implementation shifts the channel state by $-\pi$,
keeping the tracked sum consistent.

\subsection{Validation on both sides of the claim}
\emph{Static channel.} Residual versus anchor spacing $K$: 31.8\,mrad at
$K{=}1$ (22.6\% airtime), 32.9 at $K{=}5$ (14.9\%), 34.0 at $K{=}20$
(13.5\%); the fully two-way control is 27.9\,mrad at 26.0\%. A
$20\times$ change in reciprocity cadence moves the residual by 2\,mrad:
anchor cadence has decoupled from oscillator quality.

\emph{Moving channel.} Doppler at 915\,MHz is $f_D=3.05$\,Hz per m/s;
TDL-D's dominant line-of-sight tap acquires a quasi-constant Doppler
within $\pm f_D$, so $\phi_c$ drifts up to $2\pi f_D T$ per interval ---
0.19\,rad at 0.2\,m/s, 0.48\,rad at 0.5\,m/s, against the static prior
$\sigma_c=0.01$. Sweep (60 intervals, seed 0):
\begin{center}
\begin{tabular}{cclcc}
\toprule
speed (m/s) & $f_D$ (Hz) & scheme ($K$, channel prior) & residual & gain\\
\midrule
0.2 & 0.61 & control (fully two-way) & 43\,mrad & 99.95\%\\
0.2 & 0.61 & hybrid $K{=}1$, static prior & 68\,mrad & 99.88\%\\
0.2 & 0.61 & hybrid $K{=}5$, static prior & 446\,mrad & 95.3\%\\
0.2 & 0.61 & hybrid $K{=}5$, matched prior & 235\,mrad & 98.7\%\\
0.2 & 0.61 & hybrid $K{=}20$, static prior & 1752\,mrad & 55.1\%\\
0.5 & 1.52 & control (fully two-way) & 33\,mrad & 99.97\%\\
0.5 & 1.52 & hybrid $K{=}1$, static prior & 128\,mrad & 99.6\%\\
0.5 & 1.52 & hybrid $K{=}1$, matched prior & 44\,mrad & 99.95\%\\
0.5 & 1.52 & hybrid $K{=}5$, matched prior & 208\,mrad & 98.9\%\\
0.5 & 1.52 & hybrid $K{=}20$, matched prior & 530\,mrad & 93.4\%\\
\bottomrule
\end{tabular}
\end{center}
Readings. The control is immune: its measurements are channel-free at
every substep. Under the static prior, the filter attributes the channel
drift to the oscillator (larger $q_\theta$ wins the innovation) and
chases the channel with NCO corrections; degradation is monotone in both
Doppler and $K$. The matched prior $\sigma_c\approx2\pi f_DT$ recovers
fully at $K{=}1$ (44 vs.\ 33\,mrad control) and partially at sparse $K$:
correctly so, because the prior changes attribution but cannot supply
the separating information only anchors carry. Together the two
experiments establish the design rule: \textbf{reciprocity cadence is
set by channel coherence time, not oscillator quality} --- on static
links, anchors at 1\,Hz or slower suffice and total sync airtime falls
to 13--15\%; at walking-speed Doppler, anchors are needed at the
interval rate. A known residual defect: the one-way frequency
observation measures $\omega$ plus the channel Doppler and is therefore
biased at nonzero speed; a fourth state (channel Doppler rate) would
absorb it and is the next refinement.

\section{Application split: CSI-aided versus open-loop coherence}
\label{sec:csi}
The identifiability analysis of Section~\ref{sec:measure} divides
applications of these loops into two classes, and the simulator
quantifies both.

\emph{CSI-aided (closed-loop) joint transmission.} When a user feeds
back channel estimates, the two-station combining gain at the user
depends only on the \emph{differential} station phase relative to its
value at the last CSI epoch: all static biases, including the one-way
channel bias, cancel. Evaluating the one-way loop's closed-loop
oscillator residual as the differential-phase process, with pilot
estimation noise added per epoch, gives mean JT gain versus CSI refresh
cadence: 100.00\% (refresh every interval), 99.88\% (2), 99.79\% (5),
99.72\% (10), 99.56\% (20 intervals $=1$\,s). One-way synchronization
plus ordinary CSI is sufficient for cell-free-style coherent service,
with CSI staleness to 1\,s costing under 0.5\% of gain.

\emph{Open-loop coherence.} For passive coherent sensing there is no
feedback path; the carriers must align at the antennas, which requires
the two-way or hybrid loops. Their bounded residual has a consequence
specific to sensing: a \emph{closed-loop} bounded phase error (28--82
mrad, never diverging) costs a fixed 0.1--0.7\% of coherent gain
regardless of integration time, whereas open-loop (uncorrected)
oscillators decorrelate and cap the usable dwell. Integration time for
weak-target detection becomes limited by target and channel dynamics,
not by the stations' clocks.

\section{Summary and open items}
\begin{enumerate}
\item Grading discipline: estimator errors must be computed against an
independent ground truth (the oracle path), not against the estimator's
inputs; the difference was a factor of $\sim$100 here.
\item The closed-loop floor \eqref{eq:budget} is predicted by two
hand-computable terms (walk $\sqrt{kT}$, latency
$\sigma_{\hat\omega}\ell T$) and verified by ablation to within a few
mrad. The latency term has no counterpart in the published literature
through Dec.\ 2025 (19-source verified search).
\item The reference paper's algorithms reproduce at statistics level
(residuals vs.\ its Eq.\ 27) and, given reciprocity, over a physical
link, where filtered consensus and the project's EKF loop coincide at
the physical floor (82.8 vs.\ 80.8\,mrad). Applied as published, the
wrapped symmetric consensus is bistable over real channels
\eqref{eq:bistable} and locks anti-phase on realizations satisfying
$|r_0|+|\phi_c|>\pi$; verified by hand arithmetic, by a radio-free
iteration, and across seeds.
\item Cadence, not energy, moves the floor: two-tier micro-pilots reach
22.8--28\,mrad ($3\times$ the plain loop) for 7--15\% added airtime;
hybrid one-way/two-way calibration holds 33\,mrad at 13--15\% airtime on
static channels, with reciprocity cadence set by channel coherence
(validated under Doppler in both directions, including the failure of a
mismatched channel prior).
\item Open: a fourth state for channel Doppler; $N>2$ consensus over
physical links with wrap-safe (phasor-domain) updates; a
delayed-feedback PLL literature check before claiming the latency term
in print; hardware validation --- every number above is simulation.
\end{enumerate}

\section*{Sources}
\begin{itemize}
\item M.~Rashid, J.~A.~Nanzer, IEEE Trans.\ Wireless Commun.\
22(4):2789--2802, Apr.\ 2023, DOI 10.1109/TWC.2022.3213788 (PDF in this
directory; arXiv:2201.08931).
\item S.~R.~Mghabghab, A.~Schlegel, J.~A.~Nanzer, ``Impact of VCO and PLL
Phase Noise on Distributed Beamforming Arrays with Periodic
Synchronization,'' IEEE Access, 2021 (open-loop drift vs.\ update
interval).
\item Repository documents: \code{LITERATURE\_REVIEW.md} (verified
prior-art pass and novelty assessment); \code{paper\_review\_rashid\_%
nanzer.pdf}; \code{simulation\_design.pdf}; \code{README.md};
\code{hybrid\_calibration/README.md}. All numbers in this document
regenerate from \code{simulation.py} (models \code{sdr}, \code{twoway},
\code{micro}, \code{hybrid}, \code{dfpc}, \code{kfdfpc},
\code{compare}) and \code{pytest} (18 tests).
\end{itemize}

\end{document}
